\documentclass[11pt,a4paper]{article}

% ── Packages ────────────────────────────────────────────
\usepackage[utf8]{inputenc}
\usepackage[T1]{fontenc}

\usepackage[margin=2.5cm]{geometry}
\usepackage{amsmath,amssymb}
\usepackage{graphicx}
\usepackage{booktabs}
\usepackage{array}
\usepackage{longtable}
\usepackage{multirow}
\usepackage{hyperref}
\usepackage{xcolor}

\usepackage{parskip}
\usepackage{setspace}
\usepackage{fancyhdr}
\usepackage{natbib}
\usepackage{caption}
\usepackage{subcaption}
\usepackage{listings}
\usepackage{titlesec}
\usepackage{abstract}
\usepackage{authblk}
\usepackage{lineno}

% ── Colors ──────────────────────────────────────────────
\definecolor{darkblue}{RGB}{26,58,92}
\definecolor{codegray}{RGB}{245,245,245}
\definecolor{codelabel}{RGB}{100,100,100}

% ── Hyperref setup ──────────────────────────────────────
\hypersetup{
    colorlinks=true,
    linkcolor=darkblue,
    citecolor=darkblue,
    urlcolor=darkblue,
    pdftitle={Axon: A Standardized Open-Source Pipeline for Brain Organoid MEA Electrophysiology},
    pdfauthor={Metin},
    pdfkeywords={brain organoid, MEA, electrophysiology, NWB, DANDI, spike sorting, STTC}
}

% ── Section formatting ───────────────────────────────────
\titleformat{\section}{\large\bfseries}{\thesection.}{0.5em}{}
\titleformat{\subsection}{\normalsize\bfseries}{\thesubsection.}{0.5em}{}
\titleformat{\subsubsection}{\normalsize\itshape}{\thesubsubsection.}{0.5em}{}

% ── Code listing ─────────────────────────────────────────
\lstset{
    backgroundcolor=\color{codegray},
    basicstyle=\ttfamily\small,
    breaklines=true,
    frame=single,
    rulecolor=\color{gray!40},
    xleftmargin=10pt,
    xrightmargin=10pt,
    aboveskip=6pt,
    belowskip=6pt,
}

% ── Header/footer ────────────────────────────────────────
\pagestyle{fancy}
\fancyhf{}
\fancyhead[L]{\small\textit{Axon: Open-Source Organoid MEA Pipeline}}
\fancyhead[R]{\small\textit{Metin, 2026}}
\fancyfoot[C]{\small\thepage}
\renewcommand{\headrulewidth}{0.4pt}

% ── Line numbers (preprint style) ────────────────────────
\linenumbers

% ── Caption style ────────────────────────────────────────
\captionsetup{font=small, labelfont=bf, skip=4pt}

% ════════════════════════════════════════════════════════
\begin{document}

% ── Title ────────────────────────────────────────────────
\begin{center}
{\LARGE\bfseries Axon: A Standardized Open-Source Pipeline for\\[4pt]
Brain Organoid MEA Electrophysiology ---\\[4pt]
Validation Across Four Public Datasets}

\vspace{12pt}
{\large Metin}\\[4pt]
{\normalsize\itshape Independent Researcher}\\[4pt]
{\small May 2026 \quad $\cdot$ \quad Preprint --- Not peer reviewed}\\[4pt]
{\small \href{https://github.com/metin/axon}{github.com/metin/axon}}
\end{center}

\vspace{6pt}
\noindent\rule{\textwidth}{0.8pt}
\vspace{4pt}

% ── Abstract ─────────────────────────────────────────────
\begin{abstract}
\noindent
Brain organoid MEA (microelectrode array) electrophysiology is a rapidly growing 
field with no established analysis standards. Different groups use different metrics, 
thresholds, and algorithms --- making cross-study comparison difficult. We present 
\textbf{Axon} (v6.3), an open-source Python pipeline for end-to-end analysis of 
NWB-formatted brain organoid MEA recordings with direct DANDI Archive integration. 
Axon computes per-unit firing rate, inter-spike interval statistics, coefficient of 
variation, Spike Time Tiling Coefficient (STTC) cross-unit synchrony matrices, 
adaptive network burst detection, and graph-theoretic network topology via a single 
command-line interface. Version 6.3 adds full \textbf{SpikeInterface} integration 
for spike sorting of raw high-density MEA data, supporting four algorithms 
(mountainsort5, spykingcircus2, tridesclous2, simple) with DC-corrected 
RMS-based channel selection (default 64 of 942 channels) and per-unit quality 
metrics (SNR, ISI violation ratio, firing rate). We validate the pipeline on four 
public DANDI datasets spanning \textbf{29 subjects}: DANDI:001603 (human and mouse 
brain organoids and ex~vivo mouse cortex, sorted units), DANDI:001132 (human 
hippocampal organoids, HD-MEA raw signal), DANDI:000774 (in~vitro MEA, sorted 
units), and DANDI:001611 (chronic dissociated cortical culture, scalability test). 
All analysis modules are dataset-agnostic. Source code is available at 
\href{https://github.com/metin/axon}{github.com/metin/axon}.

\vspace{4pt}
\noindent\textbf{Keywords:} brain organoid, MEA electrophysiology, NWB, DANDI, 
spike sorting, SpikeInterface, STTC, graph theory, open-source, organoid intelligence
\end{abstract}

\noindent\rule{\textwidth}{0.8pt}
\vspace{4pt}

% ════════════════════════════════════════════════════════
\section{Introduction}

Brain organoids derived from induced pluripotent stem cells (iPSCs) have emerged 
as a tractable three-dimensional model for human cortical development 
\citep{lancaster2013}, disease modeling \citep{trujillo2022review}, and --- most 
recently --- as a substrate for biological computation \citep{kagan2022}. Coupled 
to microelectrode arrays (MEAs), organoids generate measurable extracellular 
electrical activity, including spontaneous firing, network bursts, and coordinated 
oscillations \citep{trujillo2019}.

Despite rapid growth, no standardized analysis pipeline exists for organoid MEA 
electrophysiology data. Different groups use different spike detection thresholds, 
burst detection algorithms, and synchrony metrics --- making cross-study comparison 
difficult \citep{smirnova2023}. The Neurodata Without Borders (NWB) format 
\citep{rubel2022} and the DANDI Archive \citep{bhatt2023} provide a common data 
standard, but the analytical layer remains fragmented.

The availability of public organoid MEA datasets --- particularly DANDI:001603 
\citep{vandermolen2026} --- creates an opportunity to develop and validate a 
standardized pipeline across multiple preparations, species, and recording 
conditions. However, these datasets contain both pre-sorted units (NWB units table) 
and raw high-density MEA recordings (ElectricalSeries only), requiring a pipeline 
that handles both data types.

We present \textbf{Axon} (v6.3), a modular Python pipeline that addresses this gap. 
Axon processes NWB files from the DANDI Archive through a standardized analysis 
chain: quality control, per-unit ISI statistics, STTC synchrony matrices, adaptive 
burst detection, and graph-theoretic network topology. For raw signal files, Axon 
integrates SpikeInterface \citep{buccino2020} for automated spike sorting with four 
algorithm options, DC-corrected RMS-based channel selection, and quality metrics per 
sorted unit. All modules are dataset-agnostic.

We validate the pipeline on 29 subjects from four public DANDI datasets, spanning 
human and mouse brain organoids, ex~vivo cortical slices, high-density hippocampal 
organoid recordings, and chronic dissociated cortical cultures. We report 
cross-dataset electrophysiological profiles, age-dependent synchrony in human 
organoids, network topology transitions, and SpikeInterface performance on HD-MEA 
data.

% ════════════════════════════════════════════════════════
\section{Methods}

\subsection{Pipeline Architecture}

Axon is implemented in Python~3.9+ and organized into 15 modules 
(Table~\ref{tab:modules}). The pipeline detects NWB file type automatically: if a 
units table is present, per-unit analysis proceeds directly; if only raw 
ElectricalSeries data is present, the user is offered threshold detection or 
SpikeInterface spike sorting. All outputs are written to a subject-specific results 
directory as CSV files and a 4-panel quicklook figure.

\begin{table}[h]
\centering
\caption{Axon v6.3 core modules.}
\label{tab:modules}
\small
\begin{tabular}{p{3.8cm} p{10.5cm}}
\toprule
\textbf{Module} & \textbf{Function} \\
\midrule
\texttt{organoid\_io.py} & NWB reading, file type detection (sorted/raw), 
RAM safety, DANDI streaming \\
\texttt{organoid\_signal.py} & Bandpass filter (300--3000~Hz), MAD-threshold 
spike detection \\
\texttt{organoid\_qc.py} & Refractory violation rate, Wilson CI, SNR \\
\texttt{organoid\_metrics.py} & ISI statistics, CV, STTC ($\Delta t = 10$~ms), 
permutation test \\
\texttt{organoid\_units\_analiz.py} & Per-unit orchestrator, STTC matrix, 
burst detection, quicklook \\
\texttt{si\_sorting.py} & SpikeInterface: 4 algorithms, RMS channel 
selection, quality metrics \\
\texttt{graph\_analiz.py} & Clustering coefficient, path length, 
small-world index $\sigma$ \\
\texttt{gruplu\_analiz.py} & Auto-grouping by prefix/duration, Mann--Whitney U \\
\texttt{organoid\_cli.py} & CLI: search, download, analyze in one command \\
\bottomrule
\end{tabular}
\end{table}

\subsection{Key Algorithms}

\subsubsection{Spike Time Tiling Coefficient (STTC)}

Cross-unit pairwise synchrony was quantified using the STTC metric 
\citep{cutts2014}, which is independent of firing rate. For two spike trains 
$A$ and $B$ with coincidence window $\Delta t = 10$~ms:

\begin{equation}
\text{STTC} = \frac{1}{2}\left[\frac{P_A - T_B}{1 - P_A \cdot T_B} + 
\frac{P_B - T_A}{1 - P_B \cdot T_A}\right]
\end{equation}

\noindent where $T_A$ is the fraction of time within $\Delta t$ of a spike in $A$, 
and $P_A$ is the fraction of $B$~spikes within $\Delta t$ of an $A$~spike. STTC 
ranges from $-1$ (anti-correlated) to $+1$ (perfectly synchronous). Matrices were 
computed for the top~20 most active units per subject.

\subsubsection{Adaptive Burst Detection}

Population firing rate was binned at 50~ms. Burst threshold was set adaptively 
using the Median Absolute Deviation (MAD):

\begin{equation}
\theta = \text{median}(r) + 3 \times \text{MAD}(r), \quad 
\text{with} \geq 3 \text{ co-firing units required}
\end{equation}

\noindent This adaptive approach avoids the fixed-threshold problem when comparing 
preparations with 10-fold differences in baseline firing rate.

\subsubsection{Graph-Theoretic Network Topology}

STTC values exceeding 0.1 defined network edges. Clustering coefficient~$C$, mean 
path length~$L$, and small-world index $\sigma = (C/C_{\text{rand}})/(L/L_{\text{rand}})$ 
were computed via NetworkX. $\sigma > 1$ indicates small-world topology.

\subsection{SpikeInterface Integration}

For raw electrophysiology data, Axon integrates SpikeInterface (v0.104.3) 
\citep{buccino2020} via the \texttt{si\_sorting.py} module. The workflow:

\begin{enumerate}
\item NWB reading with automatic ElectricalSeries path detection
\item Unsigned-to-signed dtype conversion (required for uint16 HD-MEA data)
\item \textbf{DC-offset removal} followed by per-channel RMS computation 
      (10-second window)
\item Selection of $N$ most active channels by RMS (default $N = 64$)
\item Bandpass filtering (300--3000~Hz)
\item Spike sorting: mountainsort5 (scheme1), spykingcircus2, tridesclous2, or simple
\item Quality metrics: SNR, ISI violation ratio (threshold $< 10\%$), 
      firing rate (threshold $> 0.05$~Hz)
\end{enumerate}

\noindent DC-offset correction is critical for uint16 recordings: without it, all 
channels exhibit near-identical RMS ($\approx$32,000~ADU), rendering channel 
selection meaningless. After mean subtraction, RMS reflects genuine signal amplitude 
variation (13--68~$\mu$V across channels in DANDI:001132), enabling meaningful 
selection of the most active electrode subset.

\subsection{Datasets}

Four public DANDI datasets were analyzed (Table~\ref{tab:datasets}). All data were 
accessed under CC-BY-4.0 licenses. No new biological experiments were conducted.

\begin{table}[h]
\centering
\caption{Datasets analyzed in this study.}
\label{tab:datasets}
\small
\begin{tabular}{lp{4.5cm}ccc}
\toprule
\textbf{Dataset} & \textbf{Description} & \textbf{Subjects} & 
\textbf{Data type} & \textbf{Species} \\
\midrule
DANDI:001603 \citep{vandermolen2026} & Preconfigured firing sequences, 
brain organoids & 18 & Sorted units & Human, mouse \\
DANDI:001132 & Multimodal HD-MEA, hippocampal organoids & 5 & 
Raw signal & Human \\
DANDI:000774 & Spatial/temporal propagation, in~vitro MEA & 5 & 
Sorted units & Mouse \\
DANDI:001611 & Chronic dissociated cortical cultures & 1 & 
Sorted units & Rat \\
\bottomrule
\end{tabular}
\end{table}

% ════════════════════════════════════════════════════════
\section{Results}

\subsection{Pipeline Validation --- Synthetic Tests}

Axon passed 15/15 synthetic validation tests covering: spike detection accuracy; 
STTC correctness on known inputs (identical trains: STTC\,=\,1.000; independent 
Poisson trains: STTC\,=\,$-0.015$); MAD robustness (5\% outlier injection: 
std $+403\%$, MAD $+0.8\%$); RAM overflow protection; and Wilson CI correctness 
at zero refractory violations.

\subsection{DANDI:001603 --- Cross-Category Electrophysiological Profiles}

The complete DANDI:001603 dataset (18 subjects with sorted units) was analyzed. 
Three biological categories were identified from subject metadata 
(Table~\ref{tab:001603}).

\begin{table}[h]
\centering
\caption{DANDI:001603 three-category descriptive statistics (mean\,$\pm$\,SD). 
Recording durations and ages were not controlled across categories.}
\label{tab:001603}
\small
\begin{tabular}{lcccccc}
\toprule
\textbf{Category} & $n$ & \textbf{Median Hz} & \textbf{Median CV} & 
\textbf{Median STTC} & \textbf{Burst/min} & \textbf{Active} \\
\midrule
Human organoid & 8 & $0.42 \pm 0.18$ & $1.92 \pm 0.85$ & 
$0.159 \pm 0.150$ & $9.84 \pm 6.18$ & 99\% \\
Mouse organoid & 5 & $4.42 \pm 1.48$ & $1.15 \pm 1.04$ & 
$0.014 \pm 0.030$ & $5.40 \pm 10.80$ & 49\% \\
Ex~vivo mouse cortex & 5 & $0.39 \pm 0.18$ & $2.83 \pm 1.04$ & 
$0.144 \pm 0.060$ & $10.10 \pm 5.43$ & 100\% \\
\bottomrule
\end{tabular}
\end{table}

Mouse organoids showed approximately 10-fold higher median firing rate than human 
organoids (4.42 vs.\ 0.42~Hz). Ex~vivo mouse cortex showed a firing rate profile 
comparable to human organoids (0.39~Hz) despite the same species origin, suggesting 
that the firing rate difference reflects preparation-specific factors rather than 
species genetics alone. Age differences, culture protocols, and MEA geometry were 
not controlled and remain as confounds for future investigation.

\subsection{Age-Dependent Synchrony in Human Organoids}

Human organoids were divided into two age groups based on available metadata: 
G1 (6--7~months, HO1--HO4, $n=4$) and G2 ($\approx$3.3~months/P100D, HO5--HO8, 
$n=4$). Mann--Whitney U test revealed significant differences in both median STTC 
($0.295 \pm 0.096$ vs.\ $0.023 \pm 0.008$, $p = 0.029$) and median CV 
($1.17 \pm 0.16$ vs.\ $2.68 \pm 0.53$, $p = 0.029$). Burst rate and median~Hz 
did not differ significantly ($p = 0.200$ for both).

Older human organoids thus exhibited higher cross-unit synchrony and more regular 
firing --- consistent with the maturation trajectory reported by 
\citet{trujillo2019}. However, G1 recordings lasted 3~min and G2 recordings 10~min, 
creating an age--duration confound that cannot be fully decoupled at $n=4$ per 
group. Replication with $n \geq 8$ and controlled recording duration is required.

\subsection{Network Topology --- Small-World Structure}

Graph-theoretic analysis (STTC threshold 0.1, top~20 units) revealed heterogeneous 
topology across human organoids. Two subjects (HO6 and HO7, both P100D) exhibited 
strong small-world topology ($\sigma \approx 4.3$), while older organoids 
($>$6~months) showed dense near-fully-connected networks ($\sigma \approx 1.0$). 
This suggests a topological transition during organoid maturation from sparse 
small-world to hub-dominated architecture, consistent with developmental wiring 
principles \citep{watts1998}.

\begin{table}[h]
\centering
\caption{Graph theory metrics for human organoids (DANDI:001603).}
\label{tab:graph}
\small
\begin{tabular}{lccccc}
\toprule
\textbf{Subject} & \textbf{Age} & \textbf{Density} & 
\textbf{Clustering} & \boldmath$\sigma$ & \textbf{Small-world} \\
\midrule
HO1--HO4 & P6M--P7M & 0.65--0.96 & 0.81--0.97 & 1.00--1.20 & Weak \\
HO5 & P100D & 0.068 & 0.367 & 1.000 & No \\
HO6 & P100D & 0.153 & 0.364 & 4.363 & \textbf{Strong} \\
HO7 & P100D & 0.126 & 0.390 & 4.333 & \textbf{Strong} \\
HO8 & P100D & 0.042 & 0.133 & 1.000 & No \\
\bottomrule
\end{tabular}
\end{table}

\subsection{DANDI:001132 --- SpikeInterface on HD-MEA Raw Data}

Five subjects from DANDI:001132 (\textit{Homo sapiens}, human hippocampal 
organoids) were analyzed. All files contained raw electrophysiology with 
745--972 channels at 20,000~Hz. DC-corrected RMS channel selection (64/745--972 
channels, RMS range 13--68~$\mu$V) followed by mountainsort5 (scheme1) 
sorting was applied (Table~\ref{tab:001132}).

\begin{table}[h]
\centering
\caption{DANDI:001132 SpikeInterface sorting results. 
Qualification: SNR\,$>$\,3.0, ISI violation\,$<$\,10\%, Hz\,$>$\,0.05.}
\label{tab:001132}
\small
\begin{tabular}{lccccccc}
\toprule
\textbf{Subject} & \textbf{Channels} & \textbf{Duration (s)} & 
\textbf{Units} & \textbf{Qualified} & \textbf{Median Hz} & 
\textbf{Max STTC} & \textbf{Bursts} \\
\midrule
3C-1 & 972$\to$64 & 41 & 15 & 9/15 & 1.76 & 0.196 & 3 \\
4C-1$^{*}$ & 745$\to$64 & 180 & 2 & 0/2 & 19.76 & $-0.014$ & 0 \\
5C-1 & 942$\to$64 & 51 & 6 & 4/6 & 1.11 & 0.058 & 0 \\
6C-1 & 947$\to$64 & 101 & 3 & 2/3 & 0.74 & 0.027 & 0 \\
7D-7 & 907$\to$64 & 81 & 4 & 3/4 & 1.10 & 0.028 & 0 \\
\bottomrule
\multicolumn{8}{l}{\small $^{*}$Excluded from comparative analyses (likely failed sorting; 
see Section~\ref{sec:limitations}).}
\end{tabular}
\end{table}

\subsection{DANDI:000774 --- In~Vitro MEA Validation}

Five subjects from DANDI:000774 (sorted units, mouse in~vitro MEA) were analyzed 
with per-unit analysis capped at 50~units. Recordings ranged from 3,261 to 
10,926~s with high burst counts (514--1,708). Median firing rates ranged from 
0.41 to 1.58~Hz with median STTC values 0.046--0.130, indicating moderate 
synchrony consistent with in~vitro cortical network dynamics.

\subsection{DANDI:001611 --- Chronic Culture Scalability Test}

One subject from DANDI:001611 (sub-24846, dissociated rat cortical culture, 
950 sorted units, 261,569~s recording, 974,562 total spikes) was analyzed as 
a pipeline scalability test. Axon completed the full pipeline without modification 
in 16,225~s ($\approx$4.5~hours), demonstrating that the pipeline handles 
large-scale chronic recordings. This subject is not included in organoid 
comparisons.

% ════════════════════════════════════════════════════════
\section{Discussion}

Axon provides the first open-source pipeline specifically designed for NWB-formatted 
brain organoid MEA data with DANDI integration. Validation across four datasets 
demonstrated consistent performance across sorted-units and raw-signal data types, 
multiple species and preparation types, and recording durations from 41~s to 3~days.

The age-dependent synchrony finding ($p = 0.029$, $n = 4$ per group) is consistent 
with prior reports of emergent oscillations in maturing organoids \citep{trujillo2019}. 
The strong small-world topology in HO6 and HO7 ($\sigma \approx 4.3$, both P100D) 
but not in HO5 and HO8 (also P100D, $\sigma = 1.0$) suggests that small-world 
organization may be a transient developmental feature. Age-matched longitudinal 
recordings with controlled duration are required for definitive interpretation.

The DC-offset correction step proved critical for HD-MEA channel selection. Without 
it, uint16-to-int16 converted recordings exhibit near-identical RMS ($\approx$32,000 
ADU), rendering channel selection uninformative. After mean subtraction, RMS 
variation (13--68~$\mu$V) reflects genuine differences in electrode activity.

\subsection{Limitations}
\label{sec:limitations}

\begin{itemize}
\item \textbf{Age--duration confound (DANDI:001603):} G1 (3~min) and G2 (10~min) 
recordings differ in both age and duration. Replication at $n \geq 8$ with 
controlled duration is required.
\item \textbf{Unit cap at 50 for per-unit analysis:} Prevents computational overload 
on high-unit-count preparations. Future work should implement vectorized ISI 
computation.
\item \textbf{SpikeInterface run-to-run variation:} Stochastic clustering in 
mountainsort5 produces $\pm$1 unit variation per run. This should be 
quantified systematically.
\item \textbf{4C-1 sorting failure:} Two units, 0 qualified, median Hz 19.76 for 
a 180-s recording --- anomalous. Likely reflects unresolved multi-unit activity.
\item \textbf{No isolation distance or L-ratio:} Waveform-based quality metrics 
beyond SNR are not yet computed automatically.
\end{itemize}

% ════════════════════════════════════════════════════════
\section{Conclusions}

Axon provides a standardized, reproducible, open-source analytical baseline for 
brain organoid MEA electrophysiology. By integrating DANDI access, SpikeInterface 
spike sorting, and graph-theoretic analysis in a single command-line tool, Axon 
reduces the analytical barrier for groups generating or analyzing organoid MEA data. 
Validation across 29 subjects from four public datasets demonstrates consistent 
performance across diverse recording conditions and preparation types.

The pipeline is designed as a computational companion to organoid-oi v3 
\citep{metin2026oi}, a closed-loop organoid intelligence simulation system: Axon 
quantifies organoid dynamics from existing recordings, while organoid-oi provides 
the computational framework for future closed-loop hardware experiments.

% ════════════════════════════════════════════════════════
\section*{Data and Code Availability}

All datasets are publicly available on the DANDI Archive under CC-BY-4.0:
DANDI:001603, DANDI:001132, DANDI:000774, DANDI:001611 
(\href{https://dandiarchive.org}{dandiarchive.org}).

Axon source code: \href{https://github.com/metin/axon}{github.com/metin/axon}.

organoid-oi v3: 
\href{https://github.com/metin/organoid-oi}{github.com/metin/organoid-oi}.

No new biological data were generated for this study.

% ════════════════════════════════════════════════════════
\bibliographystyle{plainnat}
\bibliography{references}

\end{document}
